% theme: definition_and_values_of_trigonometric_functions
\MajorQuestionLesson{三角関数の定義と値}
\SetRowSpace{3mm}

\TypeQuestionDouble{角 $\theta$ の動径上に、原点 $ {\rm O}$ と異なる点 $ {\rm P}(x,y)$ をとり、$r=\Lseg{OP}$ とします。次の値を $x,\ y,\ r$ を用いて表しなさい。}{teigi}
\QQRowThree{$\sin\theta$}{}{$\tan\theta$}{}{$\cos\theta$}{}

\TypeQuestionSingle{角 $\theta$ の動径上に次の点 $ {\rm P}$ があるとき、$r=\Lseg{OP}$ の値を求めなさい。}{hankei}
\QQRowThree{$ {\rm P}(6,8)$}{}{$ {\rm P}(-12,5)$}{}{$ {\rm P}(8,-15)$}{}
\QQRowThree{$ {\rm P}(-9,-12)$}{}{$ {\rm P}(20,-21)$}{}{$ {\rm P}(-7,24)$}{}

\TypeQuestionDouble{角 $\theta$ の動径が次の点 $ {\rm P}$ を通るとき、$\sin\theta,\ \cos\theta,\ \tan\theta$ の値を求めなさい。}{zahyou}
\QQRow{$ {\rm P}(3,4)$}{}{$ {\rm P}(-5,12)$}{}
\QQRow{$ {\rm P}(7,-24)$}{}{$ {\rm P}(-8,-15)$}{}
\QQRow{$ {\rm P}(0,6)$}{}{$ {\rm P}(-9,0)$}{}

\TypeQuestionSingle{次の値を求めなさい。}{kihonchi}
\QQRowThree{$\sin 0$}{}{$\tan\dfrac{1}{4}\pi$}{}{$\cos\dfrac{1}{3}\pi$}{}
\QQRowThree{$\cos\dfrac{1}{2}\pi$}{}{$\sin\dfrac{1}{6}\pi$}{}{$\tan\dfrac{1}{3}\pi$}{}
\QQRowThree{$\tan 0$}{}{$\cos 0$}{}{$\sin\dfrac{1}{4}\pi$}{}
\QQRowThree{$\sin\dfrac{1}{3}\pi$}{}{$\tan\dfrac{1}{6}\pi$}{}{$\cos\dfrac{1}{6}\pi$}{}
\QQRowThree{$\cos\dfrac{1}{4}\pi$}{}{$\sin\dfrac{1}{2}\pi$}{}{$\tan\dfrac{1}{2}\pi$}{}

\TypeQuestionDouble{角 $\theta$ の動径が指定された象限にあるとき、示された三角関数の値が正か負かを答えなさい。}{fugou}
\QQRowThree{$\theta$ は第2象限、$\sin\theta$}{}{$\theta$ は第4象限、$\tan\theta$}{}{$\theta$ は第3象限、$\cos\theta$}{}
\QQRowThree{$\theta$ は第1象限、$\tan\theta$}{}{$\theta$ は第3象限、$\sin\theta$}{}{$\theta$ は第2象限、$\cos\theta$}{}
\QQRowThree{$\theta$ は第4象限、$\cos\theta$}{}{$\theta$ は第2象限、$\tan\theta$}{}{$\theta$ は第3象限、$\tan\theta$}{}

\LessonPageBreak

\TypeQuestionSingle{次の値を求めなさい。}{degree_value}
\QQRowThree{$\sin 150^\circ$}{}{$\tan 330^\circ$}{}{$\cos 225^\circ$}{}
\QQRowThree{$\cos(-120^\circ)$}{}{$\sin 750^\circ$}{}{$\tan(-585^\circ)$}{}
\QQRowThree{$\tan 405^\circ$}{}{$\cos(-690^\circ)$}{}{$\sin(-210^\circ)$}{}

\TypeQuestionSingle{次の値を求めなさい。}{radian_value}
\QQRowThree{$\sin\dfrac{2}{3}\pi$}{}{$\tan\dfrac{5}{6}\pi$}{}{$\cos\dfrac{3}{4}\pi$}{}
\QQRowThree{$\cos\dfrac{4}{3}\pi$}{}{$\sin\dfrac{7}{6}\pi$}{}{$\tan\dfrac{3}{4}\pi$}{}
\QQRowThree{$\tan\dfrac{5}{4}\pi$}{}{$\cos\dfrac{11}{6}\pi$}{}{$\sin\dfrac{5}{3}\pi$}{}
\QQRowThree{$\sin\dfrac{3}{2}\pi$}{}{$\tan\dfrac{2}{3}\pi$}{}{$\cos\pi$}{}
\QQRowThree{$\cos\dfrac{7}{6}\pi$}{}{$\sin\dfrac{3}{4}\pi$}{}{$\tan\dfrac{11}{6}\pi$}{}

\TypeQuestionSingle{次の値を求めなさい。}{general_value}
\QQRowThree{$\sin\left(-\dfrac{1}{6}\pi\right)$}{}{$\tan\dfrac{11}{3}\pi$}{}{$\cos\left(-\dfrac{3}{4}\pi\right)$}{}
\QQRowThree{$\cos\dfrac{9}{4}\pi$}{}{$\sin\dfrac{13}{6}\pi$}{}{$\tan\left(-\dfrac{2}{3}\pi\right)$}{}
\QQRowThree{$\tan\left(-\dfrac{5}{4}\pi\right)$}{}{$\cos\dfrac{17}{6}\pi$}{}{$\sin\left(-\dfrac{7}{4}\pi\right)$}{}
\QQRowThree{$\sin\dfrac{25}{6}\pi$}{}{$\tan\dfrac{19}{4}\pi$}{}{$\cos\left(-\dfrac{11}{3}\pi\right)$}{}
